% Part: counterfactuals
% Chapter: introduction
% Section: counterfactuals

\documentclass[../../../include/open-logic-section]{subfiles}

\begin{document}

\olfileid{cnt}{int}{cnt}

\olsection{شرطی‌های خلاف واقع}

یکی از صورت‌های بسیار رایج و مهم ساختارهای انگلیسی «اگر \dots، آنگاه
\dots» با صورت التزامی گذشتهٔ فعل \emph{بودن} ساخته می‌شود: «اگر چنین
می‌بود که \dots، آنگاه چنین می‌بود که \dots». چون مقدم چنین شرطی‌ای
معمولاً کاذب، یعنی خلاف واقع، است، آن‌ها را \emph{شرطی‌های خلاف واقع}
می‌نامند (و چون از صورت التزامیِ فعل \emph{بودن} استفاده می‌کنند،
\emph{شرطی‌های التزامی} نیز نامیده می‌شوند). این‌ها از شرطی‌های
\emph{اخباری} متمایزند که صورت «اگر چنین است که \dots، آنگاه چنین است
که \dots» دارند. شرطی‌های خلاف واقع و اخباری از حیث شرایط صدق متفاوت‌اند.
نمونهٔ مشهور آدامز را در نظر بگیرید:
\begin{quote}
  اگر اسوالد کندی را نکشته است، کس دیگری او را کشته است.
  
  اگر اسوالد کندی را نکشته بود، کس دیگری او را می‌کشت.
\end{quote}
اولی اخباری و دومی خلاف واقع است. اولی آشکارا صادق است: می‌دانیم
رئیس‌جمهور جان اف. کندی به دست \emph{کسی} کشته شد؛ و اگر آن شخص
(برخلاف گزارش وارن) لی هاروی اسوالد نبود، کس دیگری کندی را کشت. دومی
چیز متفاوتی می‌گوید. ادعا می‌کند اگر اسوالد کندی را نکشته بود، یعنی اگر
تیراندازی دالاس رخ نداده یا ناموفق مانده بود، تاریخ پس از آن چنان پیش
می‌رفت که سوءقصد دیگری موفق می‌شد. برای صدق آن باید نیروهایی قدرتمند
برای تضمین مرگ جان اف. کندی توطئه کرده می‌بودند (چنان‌که بسیاری از
هواداران نظریه‌های توطئه دربارهٔ کندی باور دارند).

هنوز محل مناقشه است که آیا شرطی \emph{اخباری} را شرطی مادی به‌درستی
صورت‌بندی می‌کند؛ به‌ویژه آیا می‌توان پارادوکس‌های شرطی مادی را به
شیوه‌ای «توضیح داد» که با تعیین شرایط صدق شرطی‌های اخباری انگلیسی
سازگار باشد. در مقابل، نامناقشه است که شرطی‌های خلاف واقع را نمی‌توان
با شرطی‌های مادی به‌درستی نمادگذاری کرد. دلیل روشن است: هرچند مقدم
شرطی‌های خلاف واقع عموماً کاذب است، همهٔ شرطی‌های خلاف واقع با مقدم
کاذب صادق نیستند. برای نمونه، اگر گزارش وارن را بپذیرید و توطئه‌ای برای
ترور کندی در کار نبوده باشد، شرطی خلاف واقع آدامز یک مثال است.

شرطی‌های خلاف واقع در استدلال علّی نقشی مهم دارند: نمونه‌ای اصلی از
کاربرد آن‌ها بیان روابط علّی است. مثلاً کشیدن کبریت موجب روشن‌شدن آن
می‌شود و می‌توان این امر را چنین بیان کرد: «اگر این کبریت کشیده می‌شد،
روشن می‌شد.» شرطی‌های مادی و عموماً شرطی‌های اخباری نمی‌توانند این امر
را بیان کنند: «کبریت کشیده می‌شود~$\lif$ کبریت روشن می‌شود» اگر کبریت
هرگز کشیده نشود صادق است، مستقل از اینکه در صورت کشیده‌شدن چه رخ
می‌داد. بدتر آنکه «کبریت کشیده می‌شود~$\lif$ کبریت به دسته‌گلی بدل
می‌شود» نیز اگر کبریت هرگز کشیده نشود صادق است، حال آنکه اگر کشیده
می‌شد یقیناً به دسته‌گل بدل نمی‌شد.

هنوز دربارهٔ اینکه منطق درست شرطی‌های خلاف واقع دقیقاً چیست اختلاف نظر
وجود دارد. استالنیکر و لوئیس تحلیلی اثرگذار از آن‌ها عرضه کردند. بنا بر
تحلیل ایشان، شرطی خلاف واقع «اگر چنین می‌بود که~$S$، آنگاه چنین می‌بود
که~$T$» صادق است اگر و تنها اگر~$T$ در وضعیت خلاف واقع («جهان ممکن»)
که از همه به جهان بالفعل شبیه‌تر است و~$S$ در آن صادق است، صادق باشد.
این تحلیل «وجودشناختی» نامیده می‌شود، زیرا به هستی‌شناسی جهان‌های ممکن
ارجاع می‌دهد. تحلیل‌های دیگر از احتمال‌های شرطی یا نظریه‌های بازنگری
باور استفاده می‌کنند. منطق‌های پیشنهادی گوناگونی برای شرطی‌های خلاف
واقع وجود دارد. حتی یک منطق واحد لوئیس--استالنیکر نیز وجود ندارد:
هرچند استالنیکر و لوئیس گزارش‌هایی مشابه با ارجاع به نزدیک‌ترین
جهان‌های ممکن پیشنهاد کردند، فرض‌های متفاوتشان به استنتاج‌های معتبر
متفاوتی می‌انجامد.

\end{document}
